\section{Uniform retraction to $\RES$}
\label{sec:retract}

{\normalfont\small\color{SectionFourProposal}
\noindent\textbf{Proposed main-text revisions for author review.}
The passages in dark orange in this section are proposed replacements and
proof details; they have not yet been approved. The referee annotations
retain their original status.\par}

\begin{defn}[Dimension-free $\RES$-categories]\label{df:res}
The objects of the category $\RES$ is obtained from $\RES[*]$ by forgetting the function $f$ in the pair $(U, f)$. Any definable bijection between two such objects is a morphism of $\RES$.

% BEGIN REFEREE INSERTION R055
\RefereeRemark{055}{Typo}{Agreement}{%
\textit{Original source line 1399.}\quad Change the verb ``is'' to ``are'' in the sentence beginning ``The objects of the category $\RES$.''
}
% END REFEREE INSERTION R055
\end{defn}

%\begin{nota}
%Denote the half interval $(0, \infty) \sub \Gamma$ simply by $\bm H$. We can associate two Euler characteristics $\chi_{g}$, $\chi_{b}$ with the $\Gamma$-sort, which are distinguished by $\chi_{g}(\bm H) = -1$ and $\chi_{b}(\bm H) = 0$. They agree on doubly bounded definable sets. Both of them shall be denoted simply by $\chi$ when no distinction is needed. See \cite[Section 9]{hrushovski:kazhdan:integration:vf} for details, where $\chi_g$ is denoted by $\chi$ and $\chi_b$ by $\chi_c$.
%\end{nota}

\begin{rem}
As  the  $\Gamma$-sort is an \omin-minimal group, there is  the \memph{geometric} \omin-minimal Euler characteristic  $\chi_g : \ggk \Gamma[*] \fun \Z$. This is roughly constructed as follows. {\color{SectionFourProposal}Choose a finite cell decomposition $X=\coprod_j C_j$ of a definable set $X\sub\Gamma^k$ and set
\[
\chi_g(X)=\sum_j(-1)^{\dim C_j}.
\]
The cell dimensions may vary, including dimension zero.}
This does not depend on the chosen partition of $X$ and is invariant under definable bijections. It can also be shown that $\chi_g(X \cap [-\gamma, \gamma]^k)$ stabilizes as $\gamma$ approaches $+\infty$. The induced map, denoted by $\chi_b: \ggk \Gamma[*] \fun \Z$,  is the so-called \memph{bounded} \omin-minimal Euler characteristic. Obviously the two maps $\chi_g$, $\chi_b$  coincide on doubly bounded sets. On the other hand,  $\chi_g((0,+\infty))=-1$ whereas $\chi_b((0,+\infty))=0$. Both of them shall be denoted simply by $\chi$ when no distinction is needed.

% BEGIN REFEREE INSERTION R056
\RefereeRemark{056}{Error}{Cell dimensions in the Euler characteristic}{%
\textit{Original source line 1407.}\quad The dimension of a cell need not equal the ambient $k$: a singleton in $\Gamma^k$ already contradicts the stated decomposition into open $k$-cubes. The clean correction is to take a finite cell decomposition $X=\coprod_j C_j$ and define $\chi_g(X)=\sum_j(-1)^{\dim C_j}$. If a rectilinearization description is retained, use variable dimensions $d_j\leq k$, including $d_j=0$.
}
% END REFEREE INSERTION R056

See \cite[\S~9]{hrushovski:kazhdan:integration:vf} for more details, where $\chi_g$ is denoted by $\chi$ and $\chi_b$ by $\chi_c$.
\end{rem}


\begin{nota}\label{res:note}
Let $[\A] \in \ggk \RES$ denote the class of the affine line over the residue field. The class of the multiplicative torus $[\A] - 1$ over the residue field is written as $[\G_m]$; note that the multiplicative identity $1$ of $\ggk \RES$ is indeed the class $[1]$, but $[1]$ is not the multiplicative identity of the graded ring $\ggk \RES[*]$.

{\color{SectionFourProposal}In the graded ring, $[\A]$ denotes the homogeneous degree-one class $[\G_m]+[1]\in\ggk\RES[1]$, where $[1]$ is the degree-one point class and the multiplicative identity is the degree-zero point class $1$. Every graded localization by $[\A]$ below uses this homogeneous class. After forgetting the grading, the identity is $[\A]=[\G_m]+1$.}

% BEGIN REFEREE INSERTION R057
\RefereeRemark{057}{Clarification}{Define the homogeneous affine-line class}{%
\textit{Original source line 1414.}\quad The graded and ungraded uses of $[\mathbb A]$ still need an explicit convention. In $\ggk\RES[*]$ set $\mathbb L_1=[\mathbb G_m]+[1]\in\ggk\RES[1]$, whereas in $\ggk\RES$ set $\mathbb L=[\mathbb G_m]+1$. Specify that every graded localization by $[\mathbb A]$ below means localization by $\mathbb L_1$, and that $!I[*]$ is generated by homogeneous degree-one differences. This is not merely cosmetic: $[\mathbb G_m]+1$ in the graded ring is inhomogeneous.
}
% END REFEREE INSERTION R057

Let $!I$ be the ideal of $\ggk \RES$ generated by the elements $[\gamma^\sharp] - [\G_m]$, where $\gamma \in \Gamma$ is definable. The quotient ring $\ggk \RES / !I$ is written as $\sggk \RES$. {\color{SectionFourProposal}In the graded ring, $!I[*]$ is generated by the homogeneous degree-one differences $[\gamma^\sharp]-[\G_m]\in\ggk\RES[1]$, and we set $\sggk\RES[*]=\ggk\RES[*]/!I[*]$.}

The ideal of $\ggk \mgRES[*]$ generated by the  elements
\[
[(\G_m, \alpha + \beta)] - [(\alpha^\sharp, \beta)] \in \ggk \mgRES[1],
\]
where $\alpha, \beta \in \Gamma$ are definable, is denoted by $!\mu I[*]$, and the (graded) quotient ring $\ggk \mgRES[*] / !\mu I[*]$ by $\sggk \mgRES[*]$.
\end{nota}


The quotient maps from ``$\ggk$'' to ``$\sggk$'' will all be denoted by $\iota$, and will often be omitted when it is clear in the context that the latter one is meant. We will also use the same notation for elements when passing from the former to the latter.

Note that the category $\mgRES[*]$ and the ring $\sggk \mgRES[*]$  will be dismissed after Proposition~\ref{prop:eu:retr:k} as inadequate for our purpose, and a new construction is introduced thereafter.

\subsection{Without volume forms}\label{subs:uni:retr:novol}

The discussion here more or less follows \cite[\S~2.5]{hru:loe:lef} (also see \cite[Theorem~10.5]{hrushovski:kazhdan:integration:vf}).

For $0 \leq l \leq n$, consider the maps
\[
\epsilon^{n,l}_g : \ggk \RES[n{-}l] {\color{SectionFourProposal}\otimes_{\Z}} \ggk \Gamma[l] \fun \sggk \RES[n]
\]
given by
\begin{equation}\label{kill:gamma}
x \otimes y \efun \chi_g(y) x [\G_m]^l.
\end{equation}
{\color{SectionFourProposal}The component maps are first defined on tensor products over $\Z$. Their sum respects the homogeneous balancing relations over $\ggk\Gamma^{\fin}[*]$ after passing to $\sggk\RES[*]$: a finite Gamma class can be moved between the two factors because the differences $[\gamma^\sharp]-[\G_m]$ vanish in the target. Thus the degree-$n$ sum descends through the full graded tensor product and, together with $\Psi^{-1}$, induces a homomorphism}

% BEGIN REFEREE INSERTION R058
\RefereeRemark{058}{Clarification}{Tensor products of homogeneous pieces}{%
\textit{Original source line 1442.}\quad Strictly speaking, an individual homogeneous piece such as $\ggk\RES[n-l]$ is not a module over the whole graded ring $\ggk\Gamma^{\fin}[*]$. Define the displayed component maps first on $\ggk\RES[n-l]\otimes_{\mathbb Z}\ggk\Gamma[l]$, then show that their sum descends through the homogeneous balancing relations in the full graded tensor product. The quotient $!I[*]$ gives exactly the required compatibility.
}
% END REFEREE INSERTION R058
\[
\ggk \RV[n] \fun \sggk \RES[n].
\]
Composing this with the homomorphism
\[
\epsilon^{n,m}_g : \sggk \RES[n]  \fun \sggk \RES[m] : x \efun x [\A]^{m-n}
\]
and taking the direct sum of the resulting homomorphisms over all $n$ with $n \leq m$, we obtain a homomorphism
\[
\epsilon^m_g : \ggk \RV[{\leq} m] \fun \sggk \RES[m].
\]
These homomorphisms $\epsilon^m_g$, $m = 0, 1, \ldots$, are compatible in the obvious sense. Consequently, they give rise to a homomorphism
\[
\bb E_g : \ggk \RV[*] \fun \sggk \RES[*][[\A]^{-1}]
\]
whose range is precisely the zeroth piece $(\sggk \RES[*][[\A]^{-1}])_0$ of the $\Z$-graded ring on the righthand side. We have
\begin{equation}\label{eg:isp:van}
{\color{SectionFourProposal}\bb E_g([\RV^{\ccirc}] + 1 - [1]) = - [\G_m][\A]^{-1} + 1 - [1][\A]^{-1} = 0.}
\end{equation}

% BEGIN REFEREE INSERTION R059
\RefereeRemark{059}{Error}{Missing degree-one point factor}{%
\textit{Original source line 1461.}\quad In the graded target, $\mathbb E_g([1])=[1][\mathbb A]^{-1}$, not $[\mathbb A]^{-1}$. Thus the correct computation is $-[\mathbb G_m][\mathbb A]^{-1}+1-[1][\mathbb A]^{-1}=0$, using the homogeneous identity $[\mathbb A]=[\mathbb G_m]+[1]$. As written, the final term has degree $-1$, while the first two have degree $0$. The formula $\mathbb E_g([1])=[\mathbb A]^{-1}$ becomes correct only after forgetting the grading.
}
% END REFEREE INSERTION R059
Thus $\bb E_g$ induces an eponymous homomorphism
\[
\bb E_g : \ggk \RV[*] / (\bm P - 1) \fun \sggk \RES[*][[\A]^{-1}].
\]

There is a parallel construction, where  we replace $\chi_g$ with $\chi_b$ and $[\A]^{-1}$ with $[1]^{-1}$. The latter replacement is needed so to make $(\bm P - 1)$ vanish as in (\ref{eg:isp:van}). The resulting homomorphism is
\begin{equation}\label{EB}
\bb E_b : \ggk \RV[*] / (\bm P - 1) \fun \sggk \RES[*][[1]^{-1}],
\end{equation}
whose range is precisely the zeroth piece  $(\sggk \RES[*][[1]^{-1}])_0$.

\begin{rem}\label{eb:no:eg}
The zeroth graded piece $(\sggk \RES[*][[\A]^{-1}])_0$ of $\sggk \RES[*][[\A]^{-1}]$ is canonically isomorphic to the colimit  of the groups $\sggk \RES[n]$ with the transition morphisms $\epsilon^{n,m}_g $, which is actually what appears in the construction above. So there is an injective homomorphism from $(\sggk \RES[*][[\A]^{-1}])_0$ to the similar-looking ring $\sggk \RES[[\A]^{-1}]$, induced by the assignment $x \efun x[\A]^{-n}$ for $x \in \sggk \RES[n]$. Note that this homomorphism is not surjective.

% BEGIN REFEREE INSERTION R060
\RefereeRemark{060}{Error; revised}{The asserted injection is false}{%
\textit{Original source line 1474.}\quad The natural homomorphism in this sentence need not be injective, already over an untwisted residue field $\mathbb C$. To distinguish the two constructions precisely, let $R_n$ be the Grothendieck group of complex constructible sets of dimension at most $n$, with scissor relations imposed \emph{within that category}, and set $R=\bigoplus_{n\geq0}R_n$. This is the graded $\RES$ ring in the untwisted case; $R_n$ is not by definition the $n$th dimension-filtration subgroup of the ordinary Grothendieck ring. Let $\mathbb L_1=[\mathbb A^1]_1$. Ito--Miura--Okawa--Ueda give smooth projective threefolds $X,Y$ that are not stably birational but satisfy $([X]-[Y])\mathbb L=0$ in the ordinary ring (\href{https://arxiv.org/pdf/1606.04210}{Theorems~1.1 and~1.3}). Thus $z=([X]_3-[Y]_3)\mathbb L_1^{-3}$ maps to zero in $K_0(\mathrm{Var}_{\mathbb C})[\mathbb L^{-1}]$. Yet $z\neq0$: for every $m\geq0$, the additive invariant on $R_{3+m}$ recording its top-dimensional birational classes distinguishes $X\times\mathbb A^m$ and $Y\times\mathbb A^m$. Hence no homogeneous power of $\mathbb L_1$ kills the numerator. To fit the paper's parameter conventions, take the full algebraically closed Puiseux field over $\mathbb C$ as the parameter field. Every definable value then has a named RV lift, so all $!I$ generators already vanish and $\RES[n]$ identifies with the indicated constructible-set category. Replace ``injective homomorphism'' by ``natural homomorphism.'' The comparison of $\mathbb E_g$ with $\mathbb E_b$ uses only that natural map, and the point-localized identification in the next paragraph is unaffected.
}
% END REFEREE INSERTION R060

Similarly, there is an injective homomorphism $(\sggk \RES[*][[1]^{-1}])_0 \fun \sggk \RES[[1]^{-1}] \cong \sggk \RES$, and this is indeed an isomorphism.

If the codomain of $\bb E_b$ is changed to $\sggk \RES[[\A]^{-1}]$ in the obvious way then it may be compared with $\bb E_g$. They are different since
\begin{equation}\label{eb:neq:eg:counter}
\bb E_b([1]) = 1 \dand \bb E_g([1]) = \bb E_g([\RV^{\circ\circ}]) + 1 = -[\G_m][\A]^{-1} + 1 = [\A]^{-1}.
\end{equation}
We can equalize them by forcing $[\A] = 1$ (hence $\bb E_{g}(x \otimes y) = \bb E_{b}(x \otimes y) = 0$ if $y \notin \ggk \Gamma[0]$). The (complex version of the) resulting homomorphism is indeed the one constructed in \cite[(2.5.7)]{hru:loe:lef}.
\end{rem}

The homomorphism $\bb E_b$ will be used in constructing motivic Milnor fiber below, but not $\bb E_g$. It is not that  $\bb E_g$ would not work; it is just that, conceptually, $\bb E_b$ is a better choice,  see Remarks~\ref{eb:not:eg} and \ref{rem-comp-milnor-fib}.

% BEGIN REFEREE INSERTION R168
\RefereeRemark{168}{Typo; new}{Missing article}{%
\textit{Original source line 1485.}\quad Write ``used in constructing \emph{the} motivic Milnor fiber below.''
}
% END REFEREE INSERTION R168


%\begin{rem}\label{bbE:dim}
%In retrospect, since ``dimensions'' are lost in $\sggk \RES[[\A]^{-1}]$, in computing $\bb E_g([\bm U])$ for $\bm V = (V, g) \in \RV[*]$, we may treat $x$ in (\ref{kill:gamma}) simply as an element in $\ggk \RES$, similarly for $\bb E_b$.
%\end{rem}

\subsection{With volume forms}
Let $\fn(\Gamma, \Z)$ be the set of functions $f: \Gamma \fun \Z$ such that the range of $f$ is finite and $f^{-1}(m)$ is definable for every $m \in \Z$. Addition in $\fn(\Gamma, \Z)$ is defined in the obvious way. Multiplication in $\fn(\Gamma, \Z)$ is given by the convolution product as follows. For $f, g \in \fn(\Gamma, \Z)$ and $\gamma \in \Gamma$, let $h_\gamma$ be the function $\Gamma \fun \Z$ given by $\alpha \efun f(\alpha)g(\gamma - \alpha)$. The range of $h_\gamma$ is obviously finite. For each $m \in \Z$, $h_\gamma^{-1}(m)$ is a finite disjoint union of sets of the form
\[
f^{-1}(m') \cap (\gamma - g^{-1}(m'')),
\]
where $m' m'' = m$, and hence is $\gamma$-definable. Let
\[
 h_\gamma^* = \sum_{m} m \chi(h_\gamma^{-1}(m)) \in \Z
\]
and $f * g : \Gamma \fun \Z$ be the function  given by $\gamma \efun h_\gamma^*$. It is not hard to see that, by \omin-minimality in the $\Gamma$-sort, $f * g = g * f$ and it belongs to $\fn(\Gamma, \Z)$. Thus $\fn(\Gamma, \Z)$ is a commutative ring.

% BEGIN REFEREE INSERTION R061
\RefereeRemark{061}{Clarification; revised}{Associativity of convolution}{%
\textit{Original source line 1501.}\quad Optional proof detail: add that $p_0$ is the multiplicative identity and that associativity follows by checking the additive generators in Notation~\ref{gam:equ:cons}. For example, writing $rq_\beta=\eta r$, the only nontrivial $r,q,q$ check reduces to $\eta^2=-\eta$, which holds for both $\eta=-1$ and $\eta=0$. The ring assertion is correct for both choices of $\chi$.
}
% END REFEREE INSERTION R061

\begin{nota}\label{gam:equ:cons}
In $\fn(\Gamma, \Z)$, let $r$ be the constant function $1$ and, for each definable element $\gamma \in \Gamma$, let $p_{\gamma}$,  $q_{\gamma}$ be the characteristic functions of $\gamma$, $(\gamma, \infty)$, respectively. By \omin-minimality, $\fn(\Gamma, \Z)$ is generated as a $\Z$-module by elements of the forms $r$, $p_\gamma$, $q_\gamma$. We have the following equalities:
\[
r p_\beta   = r, \quad p_\alpha p_\beta = p_{\alpha+ \beta}, \quad p_\alpha q_\beta  = q_{\alpha+ \beta}, \quad q_\alpha q_\beta = - q_{\alpha+ \beta}.
\]
In addition,
\begin{itemize}
  \item if $\chi = \chi_{g}$ then $r^2 = r q_\beta = - r$,
  \item if $\chi = \chi_{b}$ then $r^2 = r$ and $r q_\beta = 0$.
\end{itemize}
\end{nota}

\begin{rem}\label{int:Gam}
For each $\bm I = (I, \mu) \in \mG[k]$, let $\mu_I : I \fun \Gamma$ be the function given by $\gamma \efun \mu(\gamma) + \Sigma\gamma$ and $\lambda_{\bm I} : \Gamma \fun \Z$  the function given by $\gamma \efun \chi(\mu^{-1}_I(\gamma))$. By \omin-minimality, the range of $\lambda_{\bm I}$ is finite and, for every $m \in \Z$, $\lambda_{\bm I}^{-1}(m)$ is definable. So $\lambda_{\bm I} \in \fn(\Gamma, \Z)$. If $F :(I, \mu) \fun (J, \sigma)$ is a $\mG[k]$-morphism then $F$ restricts to a bijection from $\mu^{-1}_I(\gamma)$ to $\sigma_J^{-1}(\gamma)$ for every $\gamma \in \Gamma$ and hence $\lambda_{\bm I} = \lambda_{\bm J}$ and we may write $\lambda_{\bm I}$ as $\lambda_{[\bm I]}$. Let $\bm J = (J, \sigma) \in \mG[l]$. If $k=l$ then clearly $\lambda_{[\bm I]} + \lambda_{[\bm J]} = \lambda_{[\bm I]+[\bm J]}$. An easy computation shows that, for every $\gamma \in \Gamma$,
\[
(\lambda_{[\bm I]} * \lambda_{[\bm J]})(\gamma) = \chi( (\mu \sigma)^{-1}_{I \times J}(\gamma)),
\]
where $\mu \sigma : I \times J \fun \Gamma$ is the function given by $(\alpha, \beta) \efun \mu(\alpha) + \sigma(\beta)$, and hence $\lambda_{[\bm I]} * \lambda_{[\bm J]} = \lambda_{[\bm I][\bm J]}$. {\color{SectionFourProposal}Put $\Gamma_{\mathrm{def}}=\Gamma\cap\dcl^{\mathrm{eq}}(\mathbb S)$ and let $\Z[\Gamma_{\mathrm{def}}]$ be its group ring over $\Z$, viewed as the subring of $\fn(\Gamma,\Z)$ spanned by $p_\gamma$ for all $\gamma\in\Gamma_{\mathrm{def}}$. The degree-zero objects have definable constant weights, so $\ggk\mG[0]\cong\Z[\Gamma_{\mathrm{def}}]$.} It follows that the map

% BEGIN REFEREE INSERTION R062
\RefereeRemark{062}{Error}{The group ring uses definable value-group constants}{%
\textit{Original source line 1520.}\quad The degree-zero coefficient group is the group of \emph{definable} elements of $\Gamma$, not necessarily $\Gamma(\mathbb S)$. With the present conventions, put $\Gamma_{\mathrm{def}}=\Gamma\cap\dcl^{\mathrm{eq}}(\mathbb S)$ and replace $\mathbb Z\Gamma$ throughout this construction by $\mathbb Z[\Gamma_{\mathrm{def}}]$; in the usual valued-field-generated setting $\Gamma_{\mathrm{def}}=\mathbb Q\otimes\Gamma(\mathbb S)$. For example, over $\mathbb C((t))$, the singleton object of $\mu\Gamma[0]$ with constant volume form $\operatorname{val}(t)/2$ is allowed and maps to $p_{1/2}$, although $\Gamma(\mathbb S)=\mathbb Z$. Thus the stated target misses legitimate degree-zero classes. This distinction is already acknowledged in the remark following Definition~\ref{def:Ga:cat}.
}
% END REFEREE INSERTION R062
\[
\lambda : \ggk \mG[*] \fun {\color{SectionFourProposal}\Z[\Gamma_{\mathrm{def}}]} \oplus X \fn(\Gamma, \Z)[X]
\]
determined by the assignment
\[
[\bm I] \efun \lambda_{[\bm I]}X^k, \quad \bm I \in \mG[k],
\]
is a graded ring homomorphism. Of course, there are two such homomorphisms $\lambda_g$ and $\lambda_b$, corresponding to the two choices $\chi = \chi_{g}$ and $\chi = \chi_{b}$.
\end{rem}


%recall Notation~\ref{comp:form}



We claim that the assignments {\color{SectionFourProposal}for $k\geq1$}
\begin{equation}\label{Gam:ass}
p_\gamma X^k \efun  [\G_m]^{k-1}[(\G_m, \gamma)], \quad q_\gamma X^k \efun  [\G_m]^{k-1}[(\{1\}, \gamma)], \quad r X^k \efun 0
\end{equation}
induce a graded ring homomorphism
\[
\psi : {\color{SectionFourProposal}\Z[\Gamma_{\mathrm{def}}]} \oplus X\fn(\Gamma, \Z)[X] \fun \sggk \mgRES[*] / ([\A]).
\]
Note that {\color{SectionFourProposal}$\Z[\Gamma_{\mathrm{def}}]$} may be regarded as a subring of $\ggk \mgRES[0]$ by identifying $p_\gamma$ with $[(\infty, \gamma)]$, which is the intended interpretation of the first assignment for $k = 0$. {\color{SectionFourProposal}The $q_\gamma X^k$ and $rX^k$ assignments apply only for $k\geq1$.} For the claim, it suffices to check that $\psi$ obeys the equational constraints for the generators above in Notation~\ref{gam:equ:cons}, which are all straightforward ({\color{SectionFourProposal}for $k,l\geq1$, the equality $\psi(q_\alpha X^k)\psi(q_\beta X^l)=\psi(-q_{\alpha+\beta}X^{k+l})$ holds since $[\A]=[\G_m]+[1]=0$ is forced}).

% BEGIN REFEREE INSERTION R063
\RefereeRemark{063}{Typo}{Restore the degrees in the multiplicativity check}{%
\textit{Original source line 1544.}\quad The symbols $q_\alpha$ alone are not degree-zero elements of the displayed domain. State the check as $\psi(q_\alpha X^k)\psi(q_\beta X^l)=\psi(-q_{\alpha+\beta}X^{k+l})$ for $k,l\geq1$. Also state explicitly that the second and third assignments in~\eqref{Gam:ass} apply only for $k\geq1$; the separate interpretation at $k=0$ applies to $p_\gamma$.
}
% END REFEREE INSERTION R063


Let $\gh$ denote the element $[(\RV^{\ccirc},  -\vrv)] \in \ggk \mgRV[1]$.

\begin{prop}\label{prop:eu:retr:k}
There are two graded ring homomorphisms
\[
  \mgE_{g}, \mgE_{b}: \ggk \mgRV[*] \fun   \sggk \mgRES[*] / ([\A])
\]
such that
\begin{itemize}
  \item $\bm P \in \ggk \mgRV[1]$ vanishes under both of them,
  \item for all $x \in \ggk \mgRES[*]$, $\mgE_{g}(x) = \mgE_{b}(x) = x/([\A])$,
  \item they are distinguished by $\mgE_{g}(\gh) = {\color{SectionFourProposal}+[1]/([\A])}$ and $\mgE_{b}(\gh) = 0$.
\end{itemize}
\end{prop}

% BEGIN REFEREE INSERTION R064
\RefereeRemark{064}{Error}{Sign of the geometric retraction on h}{%
\textit{Original source line 1560.}\quad The construction gives the opposite sign in the third bullet. For $\bm I=((0,\infty),-\mathrm{id})$, the corrected weight $\mu_I$ is identically zero, so $\lambda_g([\bm I])=-p_0X$. Consequently $\mu\mathbb E_g(\mathfrak h)=\psi(-p_0X)=-[\mathbb G_m]=+[1]$ modulo $([\mathbb A])$. The bounded value is indeed zero. Replace $-[1]/([\mathbb A])$ by $+[1]/([\mathbb A])$. This confirms the sign concern in the supplied human referee report; repeating the sign printed in the cited HK theorem does not fix its incompatibility with the present definitions.
}
% END REFEREE INSERTION R064
\begin{proof}
The product map
\[
(\iota /([\A])) \times (\psi \circ \lambda) : \ggk \mgRES[*] \times  \ggk \mG[*] \fun    \sggk \mgRES[*] /([\A])
\]
is $\ggk \mG^{\fin}[*]$-bilinear (because of  the ideal $!\mu I[*]$ in the target). We have
\[
{\mu}\Psi^{-1}(\bm P) = [1]-[(0, \infty)]  \dand {\mu}\Psi^{-1}(\gh) = [((0, \infty), -\id)].
\]
The existence of the desired homomorphisms follow from a straightforward computation.

{\color{SectionFourProposal}For the stated value on $\gh$, the corrected weight on $((0,\infty),-\id)$ is identically zero. Hence
\[
\lambda_g([((0,\infty),-\id)])=-p_0X,
\qquad \mgE_g(\gh)=\psi(-p_0X)=-[\G_m]=+[1]\pmod{([\A])}.
\]
For $\chi_b$, the same open ray has Euler characteristic zero, giving $\mgE_b(\gh)=0$.}

% BEGIN REFEREE INSERTION R065
\RefereeRemark{065}{Typo}{Agreement}{%
\textit{Original source line 1570.}\quad Change the final verb ``follow'' to ``follows,'' since the subject is ``The existence.''
}
% END REFEREE INSERTION R065
\end{proof}

The discussion above is more or less an expanded version of the proof of \cite[Theorem~10.11]{hrushovski:kazhdan:integration:vf}.

\begin{rem}
Alternatively, for $\chi = \chi_{g}$, we may replace $r X^k \efun 0$ in (\ref{Gam:ass}) with $r X^k \efun [\G_m]^{k-1}[1]$, that is, sending $r X^k$ and $q_0 X^k$ to the same element.  Then $\psi$ is still a graded ring homomorphism provided that we enlarge the ideal $!\mu I[*]$ by adding the elements $[1] - [(\{1\}, \gamma)] \in \ggk \mgRES[1]$. Consequently, for instance, we can construct a ring homomorphism
\begin{equation}\label{res:forget}
\ggk \mgRV[*] \fun \sggk \mgRES[*] / ([\A]) \to^{\phi} \sggk \RES[*] / ([\A]),
\end{equation}
where the homomorphism $\phi$ is induced by the obvious forgetful functor.

% BEGIN REFEREE INSERTION R066
\RefereeRemark{066}{Clarification; revised}{Name the enlarged quotient}{%
\textit{Original source line 1580.}\quad Optional notation clarification: the preceding sentence already stipulates enlargement of $!\mu I[*]$. In the alternative construction, name the enlarged ideal $J$ and write the middle term of~\eqref{res:forget} as $\ggk\mgRES[*]/(J+([\A]))$, so that the old and enlarged quotients are visibly distinguished. Read with that stipulated enlargement, the construction is consistent. The original assignment $rX^k\mapsto0$ needs no such additional relations.
}
% END REFEREE INSERTION R066

We shall not, however, pursue this further, since taking the quotient by the ideal $([\A])$ is detrimental to our purpose here. To avoid it, we need to work with the categories of doubly bounded objects.
\end{rem}

\begin{defn}[Doubly bounded $\RV$-categories]\label{db:RV}
An object $(U, f) \in \RV[k]$ is \memph{doubly bounded} if $U$ is. Denote the full subcategory of $\RV[k]$ of doubly bounded objects by $\RV^{\db}[k]$, and similarly for the categories $\mgRV^{\db}[k]$, $\mG^{\db}[k]$.
\end{defn}




\begin{rem}\label{tensor:bd}
By Lemma~\ref{db:to:db}, if $(U, f, \omega) \in \mgRV^{\db}[k]$ then $f(U)$, $\omega(U)$ are both doubly bounded as well. For essentially the same reason, ${\mu} \Psi$ restricts to an isomorphism
\[
{\mu} \Psi^{\db}: \gsk \mgRES[*] \otimes \gsk \mG^{\db}[*] \fun \gsk \mgRV^{\db}[*].
\]
\end{rem}

\begin{nota}\label{db:generator}
{\color{SectionFourProposal}Let $\fn^{\db}(\Gamma,\Z)$ be the subgroup of $\fn(\Gamma,\Z)$ consisting of functions whose nonzero locus is doubly bounded. For each definable $\gamma\in\Gamma$, set
\[
o_\gamma=\begin{cases}
\mathbf1_{(0,\gamma)},&\gamma>0,\\
\mathbf1_{(\gamma,0)},&\gamma<0,\\
-p_0,&\gamma=0.
\end{cases}
\]
As a $\Z$-module, $\fn^{\db}(\Gamma,\Z)$ is generated by $p_\gamma,o_\gamma$. For $\gamma>0$ we have $o_\gamma=q_0-q_\gamma-p_\gamma$, and for $\gamma<0$ we have $o_\gamma=q_\gamma-q_0-p_0$. The convention $o_0=-p_0$ is a convention for functions; it does not identify $o_0$ with the characteristic function of the empty interval $(0,0)$. The convolution products are given by the following cases.}

% BEGIN REFEREE INSERTION R067
\RefereeRemark{067}{Error}{Negative interval generators and bounded support}{%
\textit{Original source line 1600.}\quad For $\gamma<0$, the interval $(0,\gamma)$ is empty, so the definition as written cannot yield the claimed generators or quotient: the relation $o_\gamma+p_0$ would even kill the unit $p_0$. Define $o_\gamma=\mathbf1_{(\gamma,0)}$ when $\gamma<0$, retaining $o_\gamma=\mathbf1_{(0,\gamma)}$ when $\gamma>0$ and $o_0=-p_0$. Also define $\fn^{\db}(\Gamma,\mathbb Z)$ explicitly as the subgroup of functions whose nonzero locus is doubly bounded. These conventions make the subsequent formulas meaningful.
}
% END REFEREE INSERTION R067
{\color{SectionFourProposal}For $\beta\ne0$, put $a=\min(\alpha,\alpha+\beta)$ and $b=\max(\alpha,\alpha+\beta)$. Then $p_\alpha o_\beta$ is the characteristic function of $(a,b)$, and
\[
p_\alpha o_\beta=\begin{cases}
o_b-o_a-p_a,&0\leq a<b,\\
o_a-o_b-p_b,&a<b\leq0,\\
o_a+o_b+p_0,&a<0<b.
\end{cases}
\]
For $\beta=0$, $p_\alpha o_0=-p_\alpha$. Also,
\[
o_\alpha o_\beta=\begin{cases}
-o_{\alpha+\beta},&\alpha,\beta\geq0\text{ or }\alpha,\beta\leq0,\\
-o_\alpha-o_\beta-p_0,&\alpha<0<\beta\text{ or }\beta<0<\alpha.
\end{cases}
\]}
\end{nota}

% BEGIN REFEREE INSERTION R068
\RefereeRemark{068}{Clarification}{Specify which interval identity applies}{%
\textit{Original source line 1605.}\quad After correcting the negative-interval convention, the alternatives can be made unambiguous as follows. Set $a=\min(\alpha,\alpha+\beta)$ and $b=\max(\alpha,\alpha+\beta)$. For $\beta\neq0$, $p_\alpha o_\beta$ is the characteristic function of $(a,b)$: it is $o_b-o_a-p_a$ if $0\leq a<b$, $o_a-o_b-p_b$ if $a<b\leq0$, and $o_a+o_b+p_0$ if $a<0<b$. For $\beta=0$ it equals $-p_\alpha$. Likewise $o_\alpha o_\beta=-o_{\alpha+\beta}$ when $\alpha,\beta\geq0$ or $\alpha,\beta\leq0$, and equals $-o_\alpha-o_\beta-p_0$ when $\alpha$ and $\beta$ have opposite strict signs.
}
% END REFEREE INSERTION R068

\begin{nota}\label{db:gam:fn}
The homomorphism
\begin{equation*}
\lambda : \ggk \mG^{\db}[*] \fun {\color{SectionFourProposal}\Z[\Gamma_{\mathrm{def}}]} \oplus X \fn^{\db}(\Gamma, \Z)[X]
\end{equation*}
is constructed as before (there is only one such homomorphism now since  $\chi_{g}$, $\chi_{b}$ agree on doubly bounded sets).
\end{nota}


\begin{nota}\label{pgammabis}
Recall that $\RV^{\ccirc}_\gamma = \rv(\MM_{\gamma} \mi 0)$ (Notation~\ref{defn:disc}). {\color{SectionFourProposal}Let $D_{\RV}$ be the set of nonnegative definable $\gamma$ for which $\gamma^\sharp$ contains an $\mathbb S$-definable point. For each $\gamma\in D_{\RV}$, let}
\[
\bm P_\gamma = [\RV^{\ccirc} \mi \RV^{\ccirc}_{\gamma}] + [\{t_\gamma\}] - [1] \in \ggk \RV^{\db}[1],
\]
where $t_\gamma \in \gamma^\sharp$ is any definable point.  It also stands for the corresponding element  in $\ggk \mRV^{\db}[1]$ (with the constant volume form $0$). Note that, although $[\{t_\gamma\}] = [1]$ in $\ggk \RV^{\db}[1]$,  $[\{t_\gamma\}] \neq [1]$ in $\ggk \mRV^{\db}[1]$  unless $\gamma = 0$.

% BEGIN REFEREE INSERTION R069
\RefereeRemark{069}{Error}{The radius must admit a definable RV lift}{%
\textit{Original source line 1621.}\quad Not every definable $\gamma$ admits the definable point $t_\gamma$ required here. Restrict the definition to nonnegative $\gamma$ for which $\gamma^\sharp$ contains an $\mathbb S$-definable point, equivalently $\gamma\in\vrv(\RV(\dcl\mathbb S))$. In the later valued-field-generated setting this is the value group of the parameter field, rather than its divisible hull. For example, over $\mathbb C((t))$, $\gamma=1/2$ is definable but there is no definable RV element of that value. Make the same restriction in the generating set of $(\bm P_\Gamma)$. The later specialization to integral radii should refer back to this restriction.
}
% END REFEREE INSERTION R069

Clearly  $\bm P_\gamma$ does not depend on the choice of $t_\gamma \in \gamma^\sharp$. The ideal of $\ggk \mRV^{\db}[*]$   generated by the elements $\bm P_\gamma$ {\color{SectionFourProposal}with $\gamma\in D_{\RV}$} is denoted by $(\bm P_\Gamma)$. {\color{SectionFourProposal}A definable value need not admit a definable RV representative, as noted after Definition~\ref{def:Ga:cat}; the same restriction $\gamma\in D_{\RV}$ is used throughout this generating set.} The images of $(\bm P_\Gamma)$ under the obvious forgetful homomorphisms
\[
\ggk  \mgRV^{\db}[*] \fun \ggk \RV[*], \quad \ggk  \mgRV^{\db}[*] \fun \ggk  \mgRV[*]
\]
are contained in $(\bm P - 1)$, $(\bm P)$, respectively,

% BEGIN REFEREE INSERTION R070
\RefereeRemark{070}{Typo}{Closing punctuation}{%
\textit{Original source line 1627.}\quad End this sentence with a full stop, not a comma.
}
% END REFEREE INSERTION R070
\end{nota}

\begin{rem}\label{db:gam:res}
Let $!P$ denote the ideal of $\fn^{\db}(\Gamma, \Z)$ generated by the elements $p_\gamma - p_0$, $o_\gamma + p_0$. By the computation in Notation~\ref{db:generator}, $\fn^{\db}(\Gamma, \Z) / !P \cong \Z$; the quotient homomorphism is denoted by $\chi$ because a simple computation shows that if $\bm I = (I, \mu) \in \mG^{\db}[*]$ then indeed $\lambda_{\bm I} / !P = \chi(I)$. Now it is straightforward to check that the assignments
\begin{equation}\label{Gam:ass:bd}
p_\gamma X^k \efun  [\G_m]^{k}, \quad o_\gamma X^k \efun  - [\G_m]^{k}
\end{equation}
induce a graded ring homomorphism
\[
\psi^{\db} : {\color{SectionFourProposal}\Z[\Gamma_{\mathrm{def}}]} \oplus X\fn^{\db}(\Gamma, \Z)[X] \fun \sggk \RES[*].
\]
In a nutshell, the  homomorphism $\psi^{\db} \circ \lambda$ is a ``forgetful'' map  given by $[(I, \mu)] \efun \chi(I)[\G_m]^{k}$, $(I, \mu) \in \mG^{\db}[k]$.
\end{rem}


The forgetful homomorphism $\phi : \ggk \mgRES[*] \fun \sggk \RES[*]$ is similar to the one in  (\ref{res:forget}).

\begin{prop}\label{prop:eu:retr:k:db}
There is a graded ring homomorphism
\[
  \mgE^{\db}: \ggk \mgRV^{\db}[*] \fun \sggk \RES[*]
\]
such that $(\bm P_\Gamma)$ vanishes and, for all $x \in \ggk \mgRES[*]$, $\mgE^{\db}(x) = \phi(x)$.
\end{prop}
\begin{proof}
The homomorphism $\phi \times (\psi^{\db} \circ \lambda)$ is still $\ggk \mG^{\fin}[*]$-bilinear. {\color{SectionFourProposal}For $\gamma=0$, the annulus in $\bm P_0$ is empty, and the singleton classes $[\{t_0\}]$ and $[1]$ agree, since their values and constant volume forms are zero. Thus $\bm P_0=0$. For $\gamma\in D_{\RV}$ with $\gamma>0$, we have
\[
({\mu} \Psi^{\db})^{-1}(\bm P_\gamma) = [(0, \gamma)] + [\{\gamma\}] + [\{t_\gamma\}] - [1].
\]
Its image under $\phi\otimes(\psi^{\db}\circ\lambda)$ is
\[
\chi((0,\gamma))[\G_m]+\chi(\{\gamma\})[\G_m]+[1]-[1]
=-[\G_m]+[\G_m]=0.
\]
This proves vanishing for every generator. The open-interval-plus-endpoint expansion is used only when $\gamma>0$.}

% BEGIN REFEREE INSERTION R071
\RefereeRemark{071}{Error}{The zero-radius case in the proof}{%
\textit{Original source line 1657.}\quad The displayed decomposition and its vanishing computation apply only for $\gamma>0$. At $\gamma=0$, $\bm P_0=0$, but the displayed right side leaves the nonzero term $[\{0\}]$ in the $\Gamma[1]$ factor, whose image is $[\mathbb G_m]$. Handle $\gamma=0$ separately, and for $\gamma>0$ write the calculation explicitly: $\chi((0,\gamma))[\mathbb G_m]+\chi(\{\gamma\})[\mathbb G_m]+[1]-[1]=-[\mathbb G_m]+[\mathbb G_m]=0$. The artificial convention $o_0=-p_0$ concerns functions; it does not change the actual class of the empty interval $(0,0)$.
}
% END REFEREE INSERTION R071
\end{proof}

The composition of $\mgE^{\db}$ with the obvious forgetful homomorphism $\sggk \RES[*] \fun \sggk \RES$ is denoted by $\bb E^{\diamond}$.  By Remark~\ref{db:gam:res}, the diagram commutes:
\begin{equation}\label{edb:eb:com}
\bfig
 \Vtriangle(0,0)/->`->`->/<400, 400>[{\ggk \mgRV^{\db}[*]}`{\ggk \RV[*]}`\sggk \RES; `\bb E^{\diamond}`\bb E_b]
\efig
\end{equation}

% BEGIN REFEREE INSERTION R072
\RefereeRemark{072}{Clarification; revised}{Identify the map labelled E\_b}{%
\textit{Original source line 1665.}\quad Optional clarification of the diagram's shorthand: the right arrow is the quotient map $\ggk\RV[*]\to\ggk\RV[*]/(\bm P-1)$, followed by~\eqref{EB} and the degree-zero identification $(\sggk\RES[*][[1]^{-1}])_0\cong\sggk\RES$ from Remark~\ref{eb:no:eg}. This is a valid composite. Its point-localized identification is unaffected by the affine-localized injectivity issue noted above.
}
% END REFEREE INSERTION R072
which may serve as an alternative and more direct (but less comprehensive) construction of $\bb E^{\diamond}$.

\begin{rem}\label{eb:not:eg}
Changing $\sggk \RES$ to $\sggk \RES[[\A]^{-1}]$ in (\ref{edb:eb:com}), we see that  $\bb E_g$ does not commute with $\bb E^{\diamond}$: by (\ref{eb:neq:eg:counter}),
\[
\bb E_g([1]) = [\A]^{-1} \quad \text{whereas} \quad \bb E^{\diamond}([1]) = \bb E_b([1]) = 1.
\]
This is but an example of a general phenomenon, see Remark~\ref{eg:eb:factor}.


We can certainly redefine $\bb E^{\diamond}$ as the composition of $\mgE^{\db}$ with  the homomorphism $\sggk \RES[*] \fun \sggk \RES[[\A]^{-1}]$ implicit in Remark~\ref{eb:no:eg} and thereby make it compatible with $\bb E_g$, but as we have mentioned above, $\bb E_b$ is ultimately a better choice; see Remark~\ref{rem-comp-milnor-fib}.
\end{rem}

